% SOURCE_AUTHORITY: sga5_fr_workpass.tex, lines 11872--11892 (LF-normalized unit).
% SOURCE_SHA256: 791F4EFFC5E02832D5D77ED03518C8156D6F07E4C8238B03545DB93D883FBB28
\begin{statement}{Corolário}
Suponhamos, além disso, que a subcategoria $\cC_o$ seja estável sob $\Ker$ e $\Coker$ e seja, portanto, uma categoria abeliana, e denotemos por $D^b(\cC)_o$ a subcategoria triangulada plena de $D^b(\cC)$ gerada pelos complexos de $D^b(\cC)$ cujos objetos de cohomologia são isomorfos a objetos de $\cC_o$. Então o funtor de inclusão
\[
D^b(\cC_o)\longrightarrow D^b(\cC)_o
\]
é uma equivalência entre ambas as categorias e, por conseguinte, obtém-se a sequência de isomorfismos
\[
K(\cC_o)\simeq K_o(\cC)\simeq K(D^b(\cC_o))\simeq K(D^b(\cC)_o).
\tag{*}
\]
\end{statement}

\emph{Observação.} Suponhamos que a subcategoria $\cC_o$ satisfaça as condições (A) e (B), e que cada objeto de $\cC$ admita uma resolução finita pela esquerda formada por objetos de $\cC_o$. Neste caso, $\cC_o$ também satisfaz a condição (C), e o funtor canônico
\[
D_o^b(\cC)\longrightarrow D^b(\cC)
\]
é, de fato, uma equivalência de categorias. Por conseguinte, tem-se o isomorfismo
\[
K_{\cC_o}(\cC)\simeq K(\cC).
\]
\footnote{cf. (SGA 6 IV 1.6).}
